\subsection{Evaluation Datasets}
\label{appx:experiments:eval_datasets}
\textbf{SWE-bench.} SWE-bench is a widely used benchmark that evaluates AI systems on their ability to resolve GitHub issues~\citep{jimenez_swe-bench_2024}.
Given a codebase along with a description of a bug or feature, the AI system is asked to modify the codebase in such a way that the issue presented in the description is resolved.
SWE-bench consists of $2294$ such task instances, collected from real world pull requests (PRs) and issues in $12$ GitHub repositories that are predominantly Python.
As discussed in Section~\ref{sec:experiments}, the Lite and Verified subsets are curated from the main SWE-bench repository with the goal of making evaluation either more efficent or more reliable.
Since evaluation on the entirety of SWE-bench is fairly costly and does not have as many comparable references, we do not evaluate \texttt{SWE-agent-LM-32B} on the entire SWE-bench test set.

\textbf{SWE-bench Multimodal.} SWE-bench Multimodal applies SWE-bench collection strategy to $12$ additional predominantly JavaScript and TypeScript GitHub repositories, where task instances are associated with issues that have visual asset(s) in them~\citep{yang_swe-bench_2024}.
The evaluation dataset consists of $510$ task instances.
While the original work evaluates vision language models (VLMs) specifically, we do not evaluate \texttt{SWE-agent-LM-32B} which, as it is based on Qwen $2.5$ Coder Instruct, does not have the ability to process images as inputs.

\begin{table}[t]
    \centering
    \begin{minipage}[b]{0.3\textwidth}
        \centering
        \begin{tabular}{ll}
\toprule
\footnotesize jqlang/jq & 9 \\
\footnotesize redis/redis & 12 \\
\scriptsize micropython/micropython & 5 \\
\footnotesize valkey-io/valkey & 4 \\
\footnotesize nlohmann/json & 1 \\
\footnotesize fmtlib/fmt & 11 \\
\midrule
\textbf{C/C++} & 42 \\
\midrule
\scriptsize prometheus/prometheus & 8 \\
\footnotesize caddyserver/caddy & 14 \\
\footnotesize gin-gonic/gin & 8 \\
\footnotesize hashicorp/terraform & 5 \\
\footnotesize gohugoio/hugo & 7 \\
\midrule
\textbf{Go} & 42 \\
\midrule
\footnotesize briannesbitt/carbon & 10 \\
\footnotesize laravel/framework & 13 \\
\scriptsize phpoffice/phpspreadsheet & 10 \\
\scriptsize php-cs-fixer/php-cs-fixer & 10 \\
\midrule
\textbf{PHP} & 43 \\
\bottomrule
        \end{tabular}
    \end{minipage}
    \hfill
    \begin{minipage}[b]{0.3\textwidth}
        \centering
        \begin{tabular}{ll}
\toprule
\footnotesize apache/druid & 5 \\
\footnotesize reactivex/rxjava & 1 \\
\footnotesize apache/lucene & 9 \\
\footnotesize projectlombok/lombok & 17 \\
\footnotesize google/gson & 9 \\
\footnotesize javaparser/javaparser & 2 \\
\midrule
\textbf{Java} & 43 \\
\midrule
\footnotesize babel/babel & 5 \\
\footnotesize mrdoob/three.js & 3 \\
\footnotesize vuejs/core & 5 \\
\footnotesize preactjs/preact & 17 \\
\footnotesize axios/axios & 6 \\
\scriptsize immutable-js/immutable-js & 2 \\
\footnotesize facebook/docusaurus & 5 \\
\midrule
\textbf{JS/TS} & 43 \\
\bottomrule
        \end{tabular}
    \end{minipage}
    \hfill
    \begin{minipage}[b]{0.3\textwidth}
        \centering
        \begin{tabular}{ll}
\toprule
\footnotesize rubocop/rubocop & 16 \\
\footnotesize jekyll/jekyll & 5 \\
\footnotesize faker-ruby/faker & 2 \\
\footnotesize fastlane/fastlane & 7 \\
\footnotesize fluent/fluentd & 12 \\
\footnotesize jordansissel/fpm & 2 \\
\midrule
\textbf{Ruby} & 44 \\
\midrule
\footnotesize tokio-rs/axum & 7 \\
\footnotesize nushell/nushell & 5 \\
\footnotesize sharkdp/bat & 8 \\
\footnotesize burntsushi/ripgrep & 2 \\
\footnotesize uutils/coreutils & 5 \\
\footnotesize tokio-rs/tokio & 9 \\
\footnotesize astral-sh/ruff & 7 \\
\midrule
\textbf{Rust} & 43 \\
\bottomrule
        \end{tabular}
    \end{minipage}
    \caption{
Number of task instances per repository and language in the SWE-bench Multilingual evaluation set.
The entire dataset includes $300$ task instances covering $9$ languages.
    }
    \label{tab:swebml_dataset}
\end{table}


\textbf{SWE-bench Multilingual.} SWE-bench Multilingual is an evaluation dataset consisting of 300 task instances that we introduce with this work.
A single author carried out SWE-bench's collection strategy for $42$ additional GitHub repositories, covering the following $9$ programming languages: JavaScript, TypeScript, C, C++, Go, Java, PHP, Ruby, and Rust. These repositories span a wide range of application domains, including web frameworks, data storage and processing tools, core utilities, and widely used libraries. A brief summary of the dataset is presented in Table~\ref{tab:swebml_dataset}.

Like SWE-bench Verified, we curate the dataset by excluding task instances deemed by a team of three authors to have ambiguous or underspecified issue text. 
Each task instance edits (meaning additions and removals) on average $48$ lines of code.
Similar to SWE-bench and \bugs{}, the median number of Fail-to-Pass tests is one.

We introduce SWE-bench Multilingual to:
\begin{enumerate}
\item Provide a benchmark to evaluate model and agent performance across a variety of programming languages and application domains. Existing agent systems often rely on Python-specific tooling, effectively overfitting to the original SWE-bench~\citep{yang_swe-bench_2024}. Although SWE-bench Multimodal addresses this to some degree, its focus on visual inputs is a confounding factor for text-only evaluation of software engineering capabilities.
\item Remain fully compatible with SWE-bench, so current users can adopt it without changing infrastructure.
\item Keep the dataset small enough to run quickly. While concurrent work like \citet{zan2025multiswebenchmultilingualbenchmarkissue} provides more task instances in multiple languages, we purposely constrain the number of task instances so that the dataset is easy to run quickly.
\end{enumerate}

In \S\ref{appx:experiments:training_analyses}, we briefly discuss how performance by existing state of the art methods for SWE-bench is markedly worse on SWE-bench Multilingual, then offer some clear directions for potential next steps to build better agentic coding models that would involve extending \bugs{}.